%___________________________
%===   Redéfinition des marges par défaut
%------------------------------------------------------
%\usepackage[textwidth=18.6cm]{geometry}%à mettre dans le preambule perso
%\pagestyle{fancy}%à mettre dans le preambule perso


\setlength\paperheight{297mm}
\setlength\paperwidth{210mm}
\setlength{\evensidemargin}{0cm}% Marge gauche sur pages paires
\setlength{\oddsidemargin}%{0cm}%
{-0.5cm}% Marge gauche sur pages impaires
\setlength{\topmargin}{-2cm}% Marge en haut
\setlength{\headsep}{0.5cm}% Entre le haut de page et le texte
\setlength{\headheight}{0.7cm}% Haut de page
\setlength{\textheight}{25.2cm}% Hauteur de la zone de texte
\setlength{\textwidth}{17cm}% Largeur de la zone de texte


% Environnement enumerate
\renewcommand{\theenumi}{\bf\textsf{\arabic{enumi}}}
\renewcommand{\labelenumi}{\bf\textsf{\theenumi.}}
\renewcommand{\theenumii}{\bf\textsf{\alph{enumii}}}
\renewcommand{\labelenumii}{\bf\textsf{\theenumii.}}
\renewcommand{\theenumiii}{\bf\textsf{\roman{enumiii}}}
\renewcommand{\labelenumiii}{\bf\textsf{\theenumiii.}}


\usetikzlibrary{shadows,trees}


%definition des couleurs
\definecolor{fondpaille}{cmyk}{0,0,0.1,0}%\pagecolor{fondpaille}
\definecolor{gris}{rgb}{0.7,0.7,0.7}
\definecolor{rouge}{rgb}{1,0,0}
\definecolor{bleu}{rgb}{0,0,1}
\definecolor{vert}{rgb}{0,1,0}
\definecolor{deficolor}{HTML}{2D9AFF}
\definecolor{backdeficolor}{HTML}{EDEDED}%{036DD0}%dégradé bleu{666666}%dégradé gris
\definecolor{theocolor}{HTML}{C10CC7}%{HTML}{036DD0}%F4404D%rouge
\definecolor{backtheocolor}{HTML}{D3D3D3}
\definecolor{methcolor}{HTML}{008800}%12BB05}
\definecolor{backmethcolor}{HTML}{FFFACD}
\definecolor{backilluscolor}{HTML}{EDEDED}
\definecolor{sectioncolor}{HTML}{C10CC7}%{B2B2B2}%vert : {HTML}{008800}%{HTML}{2D9AFF}
\definecolor{subsectioncolor}{HTML}{C10CC7}%{B2B2B2}%vert : {HTML}{008800}%{rgb}{0.5,0,0}
\definecolor{subsubsectioncolor}{HTML}{C10CC7}
\definecolor{engcolor}{HTML}{D4D7FE}
\definecolor{exocolor}{rgb}{0,0.6,0}
\definecolor{exosoltitlecolor}{rgb}{0,0.6,0}
\definecolor{titlecolor}{rgb}{1,1,1}

%commande pour enlever les couleurs avant impression
\newcommand{\nocolor}
{\pagecolor{white}
\definecolor{gris}{rgb}{0.7,0.7,0.7}
\definecolor{rouge}{rgb}{0,0,0}
\definecolor{bleu}{rgb}{0,0,0}
\definecolor{vert}{rgb}{0,0,0}
\definecolor{deficolor}{HTML}{B2B2B2}
\definecolor{backdeficolor}{HTML}{EEEEEE}%{036DD0}%dégradé bleu{666666}%dégradé gris
\definecolor{theocolor}{HTML}{B2B2B2}
\definecolor{backtheocolor}{HTML}{EEEEEE}
\definecolor{methcolor}{HTML}{B2B2B2}
\definecolor{backmethcolor}{HTML}{EEEEEE}
\definecolor{backilluscolor}{HTML}{EEEEEE}
\definecolor{sectioncolor}{HTML}{B2B2B2}
\definecolor{subsectioncolor}{HTML}{B2B2B2}
\definecolor{subsubsectioncolor}{HTML}{B2B2B2}
\definecolor{engcolor}{HTML}{EEEEEE}
\definecolor{exocolor}{HTML}{3B3838}
\definecolor{exosoltitlecolor}{rgb}{0,0,0}
\definecolor{titlecolor}{rgb}{0,0,0}
}



%%%%%%%%%%%%%%%%%%%%%%%%%%%%%%%%%%%%%%%%%%%%%%%%%%%%%%%%%%%%%%%%%%%%%%%%%%%%%%%
%Encadrés pour Propriétés, Théorème, Définitions, exemples, exercices

\usepackage{environ}%pour pouvoir utiliser la commande \NewEnviron

%___________________________
%===    Cadre arrondi coloré
%------------------------------------------------------
%
\NewEnviron{CadreColor}{
\medskip
\begin{tikzpicture}[node distance=0 cm]
\node[draw,drop shadow,color=theocolor,very thick,fill=backtheocolor,rounded corners=5pt,anchor=north west] at(0,-0.02)
{\black\parbox{\linewidth-12pt}{\BODY}};
\end{tikzpicture}
%\medskip
}


%___________________________
%===    Cadre arrondi blanc
%------------------------------------------------------
%
\NewEnviron{Cadre}{
\medskip
\begin{tikzpicture}[node distance=0 cm]
\node[draw,very thick,rounded corners=5pt,anchor=north west] at(0,-0.02)
{\black\parbox{\linewidth-12pt}{\BODY}};
\end{tikzpicture}
%\medskip
}




%___________________________
%===    Redéfinition de la commande \chapter{•}
%------------------------------------------------------
%
\makeatletter

\renewcommand{\@makechapterhead}[1]{
{\color{blue}
        \begin{flushright}
            \fontfamily{pzc}
            \bfseries\fontsize{40}{20}\selectfont{
            #1}
            
        \end{flushright}
        \vspace{24pt}
            %\rule{\linewidth}{1pt}
            }
\setlength\parindent{0mm}
\renewcommand \footrulewidth{0.2pt}%
\renewcommand \headrulewidth{0pt}%
\pagestyle{fancy}
\fancyhf{}
\pieddepage{\leftmark}{}{\thepage / \pageref{LastPage}}
}

%Chapitre avec effet miroir en dessous
%{\begin{tikzpicture}[remember picture, overlay,shift={(current page.north east)}]
%        \node at (0.1,0.1) [below left,draw=\MaCouleur, very thick,rounded corners, inner sep=10pt, inner ysep=10pt,fill=gray!20,text width=3cm, align=center]%
%            {\small\color{\MaCouleur}
%            \bfseries \bsc{\chaptername \thechapter}};
%    \end{tikzpicture}\vspace*{-2cm}
%
%    \begin{center}
%        \begin{tikzpicture}[color={\MaCouleur}]
%            \draw (0,0) -- (0.8\linewidth,0)
%                node[midway,above] {\fontfamily{pzc}\selectfont\scalefont{3}#1}
%                node[midway,above,yscale=-1, opacity=0.5] {\fontfamily{pzc}\selectfont\scalefont{3}#1};
%        \end{tikzpicture}
%    \end{center}
%}

%Chapitre encadré en couleur
%{
%\begin{tikzpicture}
%\node[fill=theocolor,rectangle,rounded corners=5pt]{%
%\begin{minipage}{\linewidth}
%\begin{center}
%\vspace*{9pt}
%\textcolor{titlecolor}{\Large \textsc{\textbf{Chapitre \thechapter \ : \ #1}}}
%\vspace*{9pt}
%\end{center}
%\end{minipage}
%};\end{tikzpicture}
%}

\makeatother


%___________________________
%===    Commandes perso
%------------------------------------------------------
%
%\Leftrightarrow
\newcommand{\Lr}{\Leftrightarrow}

%Ancienne commande chapitre
\newcommand{\chapitre}[1]{
\begin{tikzpicture}
\node[fill=sectioncolor,rectangle,rounded corners=5pt]{%
\begin{minipage}{\linewidth}
\begin{center}
\vspace*{9pt}
\textcolor{titlecolor}{\Large \textsc{\textbf{#1}}}
\vspace*{9pt}
\end{center}
\end{minipage}
};
\end{tikzpicture}
\bigskip
}


\pagecolor{white}%couleur du fond de page

\renewcommand{\Pointilles}{%
\makebox[\linewidth]{\dotfill}
}


%QRcode généré par le package qrcode
\usepackage{qrcode}


%Pour pouvoir utiliser l'environnement verbatim
\usepackage{verbatim}

%Panneau danger (nécessite le package pstricks)
\def\danger{\begingroup
\psset{unit=1ex}%
\begin{pspicture}(0,0)(3,3)
 
\pspolygon[linearc=0.2,linewidth=0.12,linecolor=red](0,0)(1.5,2.6)(3,0)
 
\psellipse*(1.5,1.33)(0.14,0.75)\pscircle*(1.5,0.3){0.15}\end{pspicture}
%
\endgroup}%

\newcommand{\cad}{c'est-à-dire }



%%%%%%%%%%%%%%%%%%%%%%%%%%%%%%%%%%%%%%%%%%%%%%%%%%%%%%%%%%%%%%%%%%%%%%%%%%%%%%%